% SEGMENT OLP-0624-S01
% Part: history
% Chapter: biographies
% Section: gerhard-gentzen

\documentclass[../../../include/open-logic-section]{subfiles}

\begin{document}

\olfileid{his}{bio}{gen}

\olsection{கெர்ஹார்ட் கென்ட்சன்}

\olphoto{gentzen-gerhard}{கெர்ஹார்ட் கென்ட்சன்}

கட்டமைப்புசார் நிறுவல் கோட்பாட்டை உருவாக்கியவர் என்றே கெர்ஹார்ட்
கென்ட்சன் முதன்மையாக அறியப்படுகிறார்; குறிப்பாக இயல்பான வருவித்தல்,
தொடரணிக் கணியம் ஆகிய !!{derivation} முறைமைகளை அவர் உருவாக்கினார்.
அவர் 1909 நவம்பர் 24 அன்று ஜெர்மனியின் கிரைஃப்ஸ்வால்டில் பிறந்தார்.
ஆயத்தப் பள்ளிக்குச் செல்வதற்கு முன் மூன்று ஆண்டுகள் வீட்டிலேயே கல்வி
கற்றார்; பள்ளியில் சேர்ந்தபோது கல்வித் தரத்தில் தம் வகுப்புத் தோழர்கள்
பலரைவிடப் பின்தங்கியிருந்தார். இருந்தபோதிலும் சிறந்த மாணவராகத் திகழ்ந்து,
கணிதத்தில் மிகுந்த திறமை காட்டினார். பலவகை ஆர்வங்கள் அவருக்கு இருந்தன;
தமது தாய்க்காகக் கவிதைகளும் பள்ளி நாடக அரங்கிற்காக நாடகங்களும் எழுதினார்.

% SEGMENT OLP-0624-S02
கென்ட்சன் கிரைஃப்ஸ்வால்ட் பல்கலைக்கழகத்தில் படிக்கத் தொடங்கி, பின்னர்
க\"{o}ட்டிங்கன், மியூனிக், பெர்லின் ஆகிய நகரங்களுக்குச் சென்றார். ஹெர்மன்
வைலின் மேற்பார்வையில் க\"{o}ட்டிங்கன் பல்கலைக்கழகத்தில் 1933-இல் முனைவர்
பட்டம் பெற்றார். (அவருடைய பணியின் பெரும்பகுதியை பால் பெர்னாய்ஸ்
மேற்பார்வையிட்டார்; ஆனால் நாசிகள் அவரைப் பல்கலைக்கழகத்திலிருந்து
நீக்கினர்.) 1934-இல் டேவிட் ஹில்பர்ட்டின் உதவியாளராகப் பணியாற்றத்
தொடங்கினார். அதே ஆண்டில் \emph{Untersuchungen \"{u}ber das logische
Schlie\ss en I--II [தருக்கரீதியான தீர்வைப் பற்றிய ஆய்வுகள் I--II]}
என்ற ஆய்வுக் கட்டுரைகளில் தொடரணிக் கணியம், இயல்பான வருவித்தல் ஆகிய
!!{derivation} முறைமைகளை உருவாக்கினார். 1936-இல் பியானோ
அடிகோள்களின் முரணின்மையை நிறுவினார்.

% SEGMENT OLP-0624-S03
நாசிகளுடனான கென்ட்சனின் உறவு சிக்கலானது. அவருடைய வழிகாட்டியான பெர்னாய்ஸ்
ஜெர்மனியை விட்டு வெளியேற வேண்டியிருந்த அதே காலத்தில், நாசிகளின் துணை
இராணுவ அமைப்பான எஸ்.ஏ.-வின் பல்கலைக்கழகப் பிரிவில் கென்ட்சன் சேர்ந்தார்.
பல ஜெர்மானியர்களைப் போலவே அவரும் நாசிக் கட்சியில் உறுப்பினராக இருந்தார்.
போரின்போது வான்படை உளவுப் பிரிவின் தொலைத்தொடர்பு அதிகாரியாகப் பணியாற்றினார்.
ஆனால் 1942-இல் நரம்புத் தளர்ச்சி ஏற்பட்டதால் பணியிலிருந்து விடுவிக்கப்பட்டார்.
நாசிக் கட்சியிடமே அவருடைய விசுவாசம் இருந்ததா, அல்லது கல்விப் பணி
வெற்றியடைவதை உறுதிப்படுத்தக் கட்சியில் சேர்ந்தாரா என்பது தெளிவாகவில்லை.

% SEGMENT OLP-0624-S04
1943-இல் பிராகிலுள்ள ஜெர்மன் பல்கலைக்கழகத்தின் கணித நிறுவனத்தில்
கென்ட்சனுக்குப் பணி வழங்கப்பட்டது; அதை அவர் ஏற்றார். ஆனால் 1945-இல்
ஜெர்மன் ஆக்கிரமிப்புக்கு எதிராகப் பிராக் மக்கள் கிளர்ந்தெழுந்தனர். சோவியத்
படைகள் நகருக்குள் வந்து பல்கலைக்கழகத்தின் பேராசிரியர்கள் அனைவரையும்
கைதுசெய்தன. நாசி அமைப்புகளில் உறுப்பினராக இருந்ததால் கென்ட்சன் கட்டாய
உழைப்பு முகாமுக்கு அழைத்துச் செல்லப்பட்டார். அங்கு சிறையில் இருந்தபோது
ஊட்டச்சத்துக் குறைவால் 1945 ஆகஸ்ட் 4 அன்று 35 வயதில் இறந்தார்.

% SEGMENT OLP-0624-S05
\begin{reading}
கென்ட்சனின் விரிவான வாழ்க்கை வரலாற்றுக்கு \citet{Menzler-Trott2007}-ஐப்
பார்க்கவும். நாசி ஆட்சியின் கீழ் கணிதவியலாளர்களைப் பற்றிய ஆர்வமூட்டும்
நூல் \citet{Segal2014}; அதில் கென்ட்சனின் வாழ்க்கை பற்றிய சுருக்கமான
குறிப்பும் உள்ளது. தருக்கரீதியான தீர்வு பற்றிய கென்ட்சனின் ஆய்வுக்
கட்டுரைகள் மூல ஜெர்மன் மொழியில் கிடைக்கின்றன
\citep{Gentzen1935a,Gentzen1935b}. அவற்றின் ஆங்கில மொழிபெயர்ப்புகள்
\citet{Gentzen1969} என்ற ஒரே தொகுப்பில் சேர்க்கப்பட்டுள்ளன; அதில் ஒரு
வாழ்க்கை வரலாற்றுக் குறிப்பும் உள்ளது.
\end{reading}

\end{document}
